% Transcription — fonds Alexandre Grothendieck, shelfmark 15, batch 5 (pages 81-100).
% Established from the facsimile archives/batches/15/batch-05.pdf (a mirror of
% https://grothendieck.umontpellier.fr/15.pdf), per the skill
% transcribe-grothendieck. The batch file carries no cover sheet: PDF page N
% is page 80+N of the folder (the Montpellier footer prints « 81 » ... « 100 »),
% and the archivists' pencil numbers bottom-left agree leaf by leaf (81 to 100
% all legible).
%
% Pass: Opus 5.5 (claude-opus-5-5), 2026-09-23 — first pass, unchecked against the pages by a human.
%
% Ten of the twenty pages are transcribed: 81, then the odd pages 83 to 99.
% Absent: page 82, blank; the even pages 84 to 100, all versos of unrelated
% typescripts — 84, 88, 90, 92 leaves of a French typescript on étale
% cohomology and descent (Prop. 4.11, Remarque 3.13, 9.4.2-9.4.4, references
% to EGA IV and to exposés V-VI), scanned upside down; 90 carries a line in his
% hand, « Ceci achève la démonstration de 9.4. », which concerns that
% typescript only; 86 an English typescript, « VI.5 », on reflexive modules
% over an integrally closed noetherian domain (Theorem VI.5); 94, 96, 98 an
% English typescript on regular elements of semisimple algebraic groups
% (pp. -7-, -5-, -4-: §2 « Some recollections », Theorems 1.5-1.7, 1.10),
% upside down; 100 an English translation, typed in violet, « Chapter III.
% Some Applications of the Criterion of Projectivity of Abstract Complete
% Varieties » (typed p. 96, margin « p. 218 »). All skipped per the project's
% rule on typed versos.
%
% What the batch is. Two runs.
%   (1) Page 81, typescript, the end of « Localisation pour les variétés
%   abéliennes », begun on page 80 (batch 4): the sentence opening the page
%   completes b) of page 80 (« le foncteur P ↦ P ⊗_R X de la catégorie de ces
%   modules vers C est pleinement fidèle, ... »). NB on karoubian categories;
%   c) Mod(R,C), a stack if C is; NB (R commutative); d) restriction to an
%   affine open U; e) « apparentés » objects (locally isomorphic over A),
%   the full subcategory of objects related to a given one is A-linear,
%   localisable, « et c'est une catégorie de Deuring, » — where the typescript
%   stops, mid-sentence, two-thirds down the page. Nothing in his hand on this
%   leaf; the corrections are typed (x-ed out) or typed in the interline.
%   (2) Pages 83-99, the rectos, in his hand, blue-black ink, not paginated by
%   him and continuous from one leaf to the next (the even pages between are
%   the typed versos): the motivic Galois group of a semi-simple motivic
%   category whose lien is commutative, and Weil q-numbers.
%   83     M a semi-simple motivic category with gerbe G, lien L commutative,
%          hence a group G of multiplicative type over Q, G° a torus; E a
%          « faithful » motive, G(Q) ↪ End_M(E); a polarisation makes
%          A = End_M(E) a Riemann algebra with φ(u)φ(u)' = ε(u)·1_E; U = Ker ε
%          lands in the « unitary group » of A, so U° is a compact torus.
%   85     B the algebra generated by G(Q) in A: B = ∏ L_i, L_i totally
%          imaginary quadratic over K_i totally real; B* = H(Q),
%          H = ∏_i R_{L_i/Q} G_m, H^u = ∏_i R_{K_i/Q} Γ_i, a « compact » torus;
%          NB: determine ε : G → G_m and j : G_m → G.
%   87     general G: the commutative quotient lien L', its group H,
%          (Ker ε)° a compact torus. New start: a category of motives with a
%          « Frobenius » f generating G (Zariski).
%   89     a) ⇔ b) ⇔ c): (T of multiplicative type, f ∈ T(Q) generating) ⇔
%          (Galois-stable subgroup M of Q̄*) ⇔ (full subcategory M' of
%          Q-vector spaces with semi-simple endomorphism, stable by ⊗, ⊕,
%          sub-objects). Data j : G_m → T (j* : M → Z), ε : T → G_m with
%          εj(λ) = λ^α, q = ε(f). a) T°/Im j compact ⇒ b) |m| = q^{j*(m)/α}
%          ⇔ c) eigenvalues of absolute value q^{i/α}. A cancelled bracket.
%   91     proof a) ⇒ c) via T°_R = U_R·S_R, f^ν = u·j(x); cancelled NB;
%          a) equivalent to a « positive » structure on M.
%   93     « Q' : clôture tot. réelle de Q »; with α = 2, a) ⇔ λλ̄ = q^{j*(λ)};
%          Cor.; Proposition (i)-(iv): conditions equivalent for λ ∈ Q̄* to be
%          a Weil q-number (conjugates of absolute value q^{i/2}; λ^n ∈ Q'(i);
%          existence of G, f, j, ε with Ker ε compact; a commutative Riemann
%          algebra). Margin NBs: the real λ are the trivial ones
%          λ = (q^{1/2})^i; λ = p^{1/4} shows λ need not lie in Q'(i).
%   95     Prop.: K a number field, L/K quadratic, p split into p, p' ⇒ some
%          x ∈ L*, xσx = 1, with divisor h(p − p'); proof (x = y²/Ny). NB: K
%          totally real, L totally imaginary, λλ̄ = p^i with p not split ⇒
%          λ = p^{i/2}·(root of unity).
%   97     K totally real, L totally imaginary quadratic over K, q = p^n: the
%          group U of p-units λ with λλ̄ = q^i; U → Z^P has finite kernel and
%          image of finite index in the divisors p − p'; U ≃ Z^Q up to a
%          finite group.
%   99     « Conditions nécessaires pour qu'un λ ∈ Q̄ soit valeur propre de
%          Frobenius » for a variety over F_q: 1) λ algebraic integer,
%          2) λ ∈ Q'(√-1), 3) λ² = q^i u, u an absolute unit outside p
%          (⇔ |λ_α| = q^{i/2} for all conjugates), 4) (struck, « conséquence
%          de 2 et 3 !!! ») v_p(λ) + v_p'(λ) = in/2; the weight i;
%          « Conjecture fondamentale » (?): the necessary conditions are
%          sufficient, margin « Th. d'existence des motifs sur F_q »; this
%          would determine the motivic Galois group of F_q. The page ends
%          mid-sentence, « Vérifions que cela donne bien », continued on
%          page 101 (batch 6).
%
% Cross-batch continuations. Page 81 opens mid-sentence, continuing page 80
% (batch 4). Page 99 ends mid-sentence, continued on page 101 (batch 6).
% Page 83 opens with a new sentence (« M catégorie motivique ... »), preceded
% by a small « c » in the upper left corner, which may be a mark of order.
%
% Notation fixed for the batch. \mathcal{M} for his script M (the motivic
% category), \mathcal{M}' for the subcategory of page 89; M (roman) for the
% subgroup of Q̄* of pages 89-93; \mathcal{G} the gerbe, \mathcal{L},
% \mathcal{L}' its lien; G, H, T groups of multiplicative type; U for his
% script U (Ker ε on page 83, the maximal compact subtorus on page 91, the
% group of p-units on page 97 — three different objects, as on the page);
% \mathbf{Q}, \mathbf{Z}, \mathbf{R}, \mathbf{C}, \mathbf{N}, \mathbf{F}_q,
% \mathbf{G}_m for his barred capitals; \mathbf{Q}' the totally real closure
% of Q, \mathbf{Q}'(i) = \mathbf{Q}'(\sqrt{-1}); \mathfrak{p}, \mathfrak{p}'
% for his gothic p. The typescript's underlined C (the stack) is
% \underline{C}. His « tt », « tot. », « ext. », « exc. p », « val. abs. »,
% « t.m. », « hyp. », « prop. » are kept.
%
% Legibility. Pages 83-99 are a fast hand: the statements and formulas carry,
% the connective prose is partly \ill{}; pages 89 (lower third, a bracket
% cancelled by three long strokes) and 93 (margins) are the densest.
% Nothing on these leaves is dated.
%
% The dating is the inventory's own for the folder, copied verbatim. Its
% group, [10-18] « Théorie des motifs », is dated [à partir de 1958-à partir
% de 1982].
\documentclass[11pt,a4paper]{article}
% Le préambule commun à toutes les transcriptions du fonds Grothendieck.
%
% Il tient en une idée : **l'incertitude se compose, elle ne se résout pas.**
% Une transcription qui lisse un mot douteux a détruit la seule chose pour
% laquelle on la faisait — un lecteur doit pouvoir savoir, à chaque endroit,
% ce qui était sur le feuillet et ce qui a été ajouté. Les six macros qui
% suivent sont donc l'appareil critique, et non de la décoration.
%
% Les mêmes commandes sont reconnues par `scripts/render.mjs`, qui produit la
% vue de lecture du site. Une macro ajoutée ici doit l'être aussi là-bas,
% faute de quoi le rendu échouera bruyamment — ce qui est voulu : un rendu
% silencieusement faux est pire qu'un rendu absent.

\ProvidesPackage{grothendieck}[2026/08/08 Transcriptions du fonds Grothendieck]

\usepackage{amsmath,amssymb,amsthm}
\usepackage{mathrsfs}
\usepackage{stmaryrd}
% Les diagrammes commutatifs sont partout dans ce fonds : c'est la langue
% même de la géométrie algébrique de Grothendieck, et une transcription qui
% les rendrait en prose ne transcrirait rien.
\usepackage{tikz-cd}
% Français par défaut — c'est la langue des feuillets — mais l'anglais est
% chargé aussi : la traduction et le résumé sont des documents anglais, et la
% typographie française (espaces avant les deux-points) y serait une faute.
% Ces documents basculent par \selectlanguage{english} en tête de corps.
\usepackage[english,french]{babel}
\usepackage{fontspec}
\usepackage{geometry}
\geometry{a4paper,margin=25mm,marginparwidth=28mm}
\usepackage{marginnote}
\usepackage{xcolor}
% ulem, not soul. `soul`'s \st and \ul are built on a character-by-character
% scan that XeLaTeX's font machinery breaks: the document fails with an
% undefined control sequence rather than degrading. ulem does the same two jobs
% and is Unicode-engine safe. `normalem` stops it hijacking \emph, which the
% transcriptions use for Grothendieck's own emphasis.
\usepackage[normalem]{ulem}
\usepackage{hyperref}
\hypersetup{colorlinks=true,linkcolor=black,urlcolor=black,pdfborder={0 0 0}}

% --- Métadonnées du lot -----------------------------------------------------
% Reprises telles quelles par le script de rendu, qui les affiche en tête de
% la vue de lecture. Elles fixent ce que le fichier prétend couvrir : sans
% elles, une transcription flotte sans point d'attache dans un fonds de seize
% mille feuillets.

\newcommand{\folder}[1]{\gdef\grothFolder{#1}}
\newcommand{\batch}[1]{\gdef\grothBatch{#1}}
\newcommand{\leaves}[2]{\gdef\grothFirst{#1}\gdef\grothLast{#2}}
\newcommand{\foldertitle}[1]{\gdef\grothFoldertitle{#1}}
\gdef\grothFolder{?}\gdef\grothBatch{?}\gdef\grothFirst{?}\gdef\grothLast{?}\gdef\grothFoldertitle{}

% --- Le résumé --------------------------------------------------------------
%
% Il ouvre la lecture modernisée, sous le titre « Résumé », et il est composé
% en petit. Ce n'est pas un ornement : c'est le seul passage du document qui
% ne soit pas des mathématiques, et un lecteur doit pouvoir le reconnaître
% avant de l'avoir lu — puis le sauter s'il connaît déjà le sujet, ou s'y
% arrêter s'il ne le connaît pas. Le filet qui le suit marque la reprise du
% corps.

\newenvironment{resume}
  {\small\setlength{\parskip}{0.45em}}
  {\par\medskip\hrule\medskip}

% --- Mots-clés ---------------------------------------------------------------
%
% En anglais, en clôture du résumé de la lecture modernisée : le vocabulaire
% moderne sous lequel chercher ce que les feuillets construisent. Le manifeste
% du site (`npm run manifest`) les extrait pour étiqueter la cote entière —
% cette ligne est la source unique des tags, il n'y a pas de fichier à côté.
\newcommand{\keywords}[1]{\par\smallskip\noindent
  {\footnotesize\textsc{Keywords} --- \textit{#1}}\par}

% --- L'appareil critique ----------------------------------------------------

% Le numéro de feuillet, en marge. C'est l'ancre qui permet de régler un
% désaccord sur une formule en regardant *un* feuillet plutôt qu'une chemise
% de six cents. Le site s'en sert aussi pour tourner les pages du fac-similé
% quand on fait défiler la transcription.
\newcommand{\leaf}[1]{%
  \marginnote{\footnotesize\color{gray!70}\textsf{#1}}%
  \label{leaf:#1}\ignorespaces}

% Illisible : on ne devine pas. Un mot inventé qui se lit comme les autres est
% le pire résultat possible.
\newcommand{\ill}{\textcolor{red!60!black}{[\,\ldots\,]}}

% Lecture douteuse : le mot est donné, et signalé comme incertain.
\newcommand{\uncertain}[1]{\uline{#1}}

% Ajout de l'éditeur — une lettre manquante, un mot sous-entendu.
\newcommand{\add}[1]{\textcolor{blue!50!black}{[#1]}}

% Biffé par Grothendieck lui-même. Ce qu'il a rayé fait partie du document :
% une rature dit souvent où la pensée a changé de direction.
\newcommand{\struck}[1]{\sout{#1}}

% Note du transcripteur, sur l'état matériel du feuillet.
\newcommand{\note}[1]{\par\smallskip{\small\color{gray!60!black}\textit{[#1]}}\par\smallskip}

% Note marginale de Grothendieck, distincte du corps du texte.
\newcommand{\marginal}[1]{\par\smallskip\noindent\hspace{1em}%
  {\small\color{gray!60!black}#1}\par\smallskip}

% --- Titre ------------------------------------------------------------------

\AtBeginDocument{%
  \begin{center}
    {\large\bfseries Fonds Alexandre Grothendieck --- cote n\textsuperscript{o}~\grothFolder}\\[2pt]
    {\itshape\grothFoldertitle}\\[4pt]
    {\small feuillets \grothFirst--\grothLast\ (lot \grothBatch)}\\[2pt]
    {\footnotesize\color{gray!60!black}Transcription établie d'après le fac-similé numérisé
     par l'Université de Montpellier.\\
     Les lectures incertaines sont soulignées, les illisibles notées [\,\ldots\,],
     les ajouts entre crochets.}
  \end{center}
  \vspace{1em}}

\endinput


\folder{15}
\batch{5}
\pages{81}{100}
\dating{1965-[vers 1977]}
\watermark{Édition de démonstration}
\foldertitle{Théorie arithmétique. Théorie de Galois des motifs (1965) : notes manuscites (s.d.), tapuscrit (s.d.).}

\begin{document}

\section{Localisation pour les variétés abéliennes}

\note{titre du tapuscrit commencé au feuillet 80 (lot 4), dont le
feuillet 81 est la fin ; la première phrase achève l'alinéa b) du
feuillet 80 (« le foncteur $P \mapsto P \otimes_R X$ »). Les corrections
sont dactylographiées (mots barrés de x, ajouts tapés dans l'interligne) ;
rien n'est de sa main sur ce feuillet}

\page{81}
de la catégorie de ces modules vers $C$ est pleinement fidèle, et son image
essentielle est formée des objets $X'$ de $C$ qui sont localement sur $S$
isomorphes à un objet de la forme $X^n$.
\note{les deux lignes « image essentielle … localement » et « sur S
isomorphes … $X^n$ » sont tapées l'une sur l'autre ; la lecture suivie est
celle-ci}

NB La conclusion de b) est satisfaite si, au lieu de supposer $C$
localisable, on la suppose stable par facteurs directs (=karoubienne), et
\struck{en fait P} si on prend des $P$ projectifs de type fini sur $R$.

c) Soit toujours $R$ une $A$-algèbre [de type fini], et soit
$\mathrm{Mod}(R,C)$ la catégorie des objets de \uncertain{$R$}\note{sic :
le sens demande « objets de $C$ »} munis d'une opération $R$-linéaire.
C'est une catégorie $A$-linéaire, et la catégorie fibrée des
$\mathrm{Mod}(\underline{C},R)_{A'}$, (pour $A'$ \struck{\ill{}} [une
$A$-algèbre]) \struck{affine d'un ouvert affine de $S = \mathrm{Spec}(A)$}
est équivalente à la catégorie fibrée des sections de $\underline{C}$
munis d'une opération de $R$. En particulier, sa restriction au
\struck{site} petit site zariskien affine de $S$ \struck{site} s'identifie à
$\mathrm{Mod}(\underline{C},R)$ ; c'est donc un champ si $\underline{C}$ en
est un, donc $\mathrm{Mod}(R,C)$ est localisable [sur $S$] si $C$ l'est.
\note{« de type fini », « une A-algèbre » et « sur S » sont tapés dans
l'interligne ; l'ordre des arguments, $\mathrm{Mod}(R,C)$ ou
$\mathrm{Mod}(\underline{C},R)$, varie ainsi dans le tapuscrit}

NB Si $R$ est commutatif, $\mathrm{Mod}(R,C)$ est même $R$-linéaire. Notre
raisonnement \emph{ne prouve pas} que \struck{C est} $\mathrm{Mod}(R,C)$ est
localisable sur $R$ i.e.\ sur $\mathrm{Spec}(R)$ (ce qui est plus fort que
de dire qu'elle est localisable sur $S = \mathrm{Spec}(A)$).

d) \struck{Deux objets $X$, $X'$ de $C$ sont dits} Si $C$ est localisable
sur $A$, alors pour $A'$ \struck{\ill{}} anneau affine d'un ouvert affine
[$U$] de $S$, $C_{A'}$ est localisable sur $A'$, et le champ associé sur
$U$ s'identifie à la restriction du champ $\underline{C}$ à $U$.

e) Deux objets de $C$ sont dits apparentés (ou $A$-apparentés) s'ils sont
localement isomorphes sur $A$. C'est évidemment une relation d'équivalence,
et la sous-catégorie pleine de $C$ formée des objets apparentés à un objet
de $C$ est une catégorie \struck{de} $A$-linéaire localisable sur $A$, et
c'est une catégorie de Deuring,
\note{le tapuscrit s'arrête ici, au milieu de la phrase, aux deux tiers
du feuillet ; le feuillet 82 est vierge}

\section{[Gerbe motivique à lien commutatif ; nombres de Weil]}

\note{titre de l'éditeur, entre crochets : la suite manuscrite des
feuillets 83 à 99 (les rectos ; les versos pairs sont des tapuscrits
étrangers) ne porte ni titre ni pagination de sa main. Elle se lit d'un
feuillet au suivant}

\page{83}
\marginal{c}
\note{petit « c » dans le coin supérieur gauche, peut-être une marque
d'ordre}

$\mathcal{M}$ catégorie motivique [semi-simple], de gerbe $\mathcal{G}$.
Supposons que le lien $\mathcal{L}$ de $\mathcal{G}$ soit
\struck{\ill{}} commutatif, donc correspond à un groupe algébrique
[\uncertain{commutatif} de type multiplicatif] $G$ sur $\mathbf{Q}$. Ainsi,
$G^{\circ}$ est un tore. Soit $E$ un motif « fidèle ». \struck{et par $F$
\ill{}} Alors $G$ opère sur $E$, considéré comme \struck{objet} [section] de
la catégorie fibrée $\mathrm{HOM}(\mathcal{G}, \uncertain{\mathrm{Proj}})$,
et en particulier $G(\mathbf{Q})$ opère sur $E$, donc
\[
  G(\mathbf{Q}) \overset{\varphi}{\hookrightarrow} \mathrm{End}_{\mathcal{M}}(E).
\]
\note{« semi-simple » est ajouté au-dessus de la ligne, écrit « ½
simple » ; « commutatif de type multiplicatif » est ajouté sous la ligne,
avec un trait de renvoi devant « $G$ ». Le second argument de
$\mathrm{HOM}$ est souligné, comme pour un champ}
D'ailleurs, \struck{les éléments de $G(\mathbf{Q})$ sur
$U = \mathrm{Ker}(G^{\circ} \to \mathbf{G}_m)$, dans $U$ opère}
Choisissons une forme de polarisation sur $E$, d'où une structure
d'algèbre de Riemann sur $A = \mathrm{End}_{\mathcal{M}}(E)$, et on a
\ill{} pour tout $u \in G(\mathbf{Q})$,
\[
  \varphi(u)\,\varphi(u)' = \varepsilon(u)\,1_{E} .
\]
[En particulier, si $u \in U = \mathrm{Ker}\,\varepsilon$, on a
$\varphi(u)\varphi(u)' = \mathrm{id}$]. Donc l'algèbre engendrée par les
$\varphi(u)$ en $A$ est \uncertain{auto-adjointe}, et $U(\mathbf{Q})$
s'envoie comme un s-groupe dans le « groupe unitaire » de $A$, formé des
$u$ tels que $uu' = 1$. \struck{Plus précisément} Or ce groupe [formé des
\ill{} d'un \ill{} vect$/\mathbf{Q}$] est un tore compact. On trouve que
$U^{\circ}$ est un tore compact.
\note{l'ajout entre crochets, à droite au-dessus de « groupe », d'une
autre encre, est en partie illisible}

\page{85}
Plus précisément, en introduisant l'algèbre engendrée par \struck{$A$}
$G(\mathbf{Q})$ dans $A$, soit $B$. On sait que
$B = \prod L_i$, chaque $L_i$ ext.\ quadratique d'un corps de nombres $K_i$
tot.\ réel ($K_i$ déterminé uniquement \struck{par $L_i$ \ill{}} [par
$L_i$] comme \ill{} tot.\ réel de $L_i$), l'involution \ill{} induite sur
$L_i$ par la conjugaison \uncertain{par rapport} à $K_i$ (i.e.\ la
conjugaison complexe ordinaire, si $L_i \subset \mathbf{C}$), la trace est
proportionnelle, sur chaque $L_i$, à la trace ordinaire.

Ainsi $B^{*}$ n'est autre que $H(\mathbf{Q})$, où
\[
  H = \prod_{i} \Bigl( \prod_{L_i/\mathbf{Q}} \mathbf{G}_{m\,L_i} \Bigr)
\]
et le groupe unitaire $B^{*u}$ s'obtient \uncertain{comme} le groupe des
\uncertain{points} rationnels de
\[
  H^{u} = \prod_{i} \Bigl( \prod_{K_i/\mathbf{Q}} \Gamma_i \Bigr)
\]
où $\Gamma_i$ est la forme tordue de $\mathbf{G}_m$ sur $K_i$ définie par
l'extension quadratique $L_i$ de $K_i$. [\struck{Ainsi} Cela montre au
\ill{} que $H^{u}$ est un tore « compact », \struck{\ill{}} (de nature
d'ailleurs assez particulière !), et aussi $U^{\circ}$ comme sous-tore de
ce dernier \ill{}.
\note{$\prod_{L_i/\mathbf{Q}}$ et $\prod_{K_i/\mathbf{Q}}$ sont sa
notation pour la restriction des scalaires à la Weil}

N.B. Il conviendrait de déterminer $\varepsilon : G \to \mathbf{G}_m$ en
termes des \uncertain{tores normes} sur les $L_i$, et
$j : \mathbf{G}_m \to G$ \ill{} $=$ \ill{}.

\page{87}
Dans le cas où $\mathcal{G}$ est quelconque, on \ill{} définir cependant
le [lien] $\mathcal{L}' = \uncertain{\mathrm{coval}}(\mathcal{L})$, qui
est commutatif, et il correspond à une gerbe quotient $\mathcal{G}'$ de
lien commutatif $\mathcal{L}'$. D'ailleurs, $\mathcal{L}'$ correspond à un
groupe algébrique [$H$], et ce qui précède prouve que
\[
  \bigl( \mathrm{Ker}(\varepsilon : H \to \mathbf{G}_m) \bigr)^{\circ}
\]
est un tore \emph{compact}.
\note{le mot devant « $(\mathcal{L})$ » se lit « coval » ; le sens
demande le quotient commutatif maximal du lien}

\note{un double trait horizontal sépare ce qui suit}

Nous \struck{envisagerons} maintenant une \struck{gerbe} catégorie de
motifs munie d'un \struck{\ill{}} Frobenius \uncertain{tensoriel} dans
\uncertain{$\mathrm{Aut}(\mathrm{id})$} qui soit « générateur ». Ainsi
$\mathcal{L}$ est commutatif et correspond à un groupe [de type]
\struck{\ill{}} multiplicatif $G$, muni d'un élément $f$ (« frobenius »),
et $G$ est engendré (au sens de Zariski) par $f$.
\struck{Dans le cas qui nous occupe, $E$ \ill{} Précédent \ill{}}
\struck{soit}.

\page{89}
\begin{enumerate}
  \item[a)] \struck{motif} [groupe de type mult.] $T$ sur $\mathbf{Q}$,
  avec $f \in T(\mathbf{Q})$ engendrant $T$ (Zar)
\end{enumerate}
\[ \Updownarrow \]
\begin{enumerate}
  \item[b)] Sous-groupe \struck{de} $M$ de $\overline{\mathbf{Q}}^{*}$,
  stable par $\mathrm{Gal}(\overline{\mathbf{Q}}/\mathbf{Q})$
\end{enumerate}
\[ \Updownarrow \]
\begin{enumerate}
  \item[c)] Sous-catégorie strictement pleine $\mathcal{M}'$ de la
  catégorie des vectoriels sur $\mathbf{Q}$ avec endomorphisme
  [semi-simple] (i.e.\ de la catégorie des $\mathbf{Q}[X]$-modules
  [semi-simples] de dim.\ finie sur $\mathbf{Q}$) stable par $\otimes$,
  $\oplus$, sous-objet.
\end{enumerate}

\struck{\ill{}} Soit donné un groupe de t.m.\ $T$, avec « générateur »
$f \in T(\mathbf{Q})$, et soit donné
$\mathbf{G}_m \xrightarrow{\;j\;} T$ [invariant par
$\mathrm{Gal}(\overline{\mathbf{Q}}/\mathbf{Q})$], qui équivaut à se donner
$j^{*} : M \to \mathbf{Z}$ \struck{\ill{} $j^{*}(f) = q$}, et on se donne
un homomorphisme $\varepsilon : T \to \mathbf{G}_m$ (\struck{tel que}
$\varepsilon j(\lambda) = \lambda^{\alpha}$) [donnée fonctorielle
\ill{}], i.e.\ $j^{*}(\varepsilon^{*}) = \alpha$, \struck{\ill{}} i.e.\ soit
donné un objet $\Lambda$ de $\mathcal{M}'$ [de rang $1$] \ill{} ;
i.e.\ \struck{Pour $\varepsilon(f) = q$ \ill{} $|\varepsilon(f)|$}.
\marginal{$q = \varepsilon^{*} = \varepsilon(f) \in M \cap \mathbf{Q}^{*}$ ;
i.e.\ $\Lambda$ de degré $\alpha$}
\note{toute cette partie est très surchargée : mots biffés, un passage
noirci, des ajouts au-dessus et au-dessous de la ligne ; la note marginale
est cerclée}

\begin{enumerate}
  \item[a)] $T^{\circ}/\mathrm{Im}\,j$ est compact.
\end{enumerate}
\[ \Downarrow \]
\begin{enumerate}
  \item[b)] [$q > 0$ et] Pour tout $m \in M$, on a
  $|m| =$ \struck{\ill{}} $q^{j^{*}(m)/\alpha}$
\end{enumerate}
\[ \Updownarrow \]
\begin{enumerate}
  \item[c)] [$q > 0$ et] Pour tout $(E, f) \in \mathcal{M}'$, si $E$ est de
  degré $i$, les valeurs propres de $f$ dans $E_{\mathbf{C}}$ sont de
  val.\ absolue $q^{i/\alpha}$.
\end{enumerate}

En effet, les valeurs propres des \struck{$(E,f) \in \mathcal{M}$} $f$,
pour $(E,f) \in \mathcal{M}'$ \ill{}, sont les $m \in M$ avec
$j^{*}(m) = i$. Donc \struck{\ill{}} c) et b) sont équivalents.

\struck{[Ces impliquent a), car $T' = T^{\circ}/\mathrm{Im}\,j$ est un tore
\ill{} ; \ill{} $f'$ \ill{} représentation \ill{} $E$, les valeurs propres
\ill{} de val.\ abs.\ $1$ \ill{} $f'$ opérant sur $E_{\mathbf{C}}$ \ill{}
$\mathrm{Gal}(E_{\mathbf{C}})$ \ill{} $T'_{\mathbf{R}}$ est un
$\mathbf{R}$-groupe \ill{} compact \ill{} $E_{\mathbf{C}}$ est de val.\ abs.\ $1$ \ill{}]}
\note{ce crochet, jusqu'au bas du feuillet, est biffé par trois longs
traits obliques et plusieurs traits horizontaux ; seuls des fragments se
lisent}

\page{91}
Montrons que a) implique c). Soit en effet \struck{\ill{}} $E$ un
vectoriel \struck{sur $\mathbf{Q}$} [« de degré $i$ »] avec $T$ opérant.
\struck{Alors les opérations de $T_{\mathbf{R}}$ opérant sur
$E_{\mathbf{R}}$, et comme \ill{} $U$ le \ill{} groupe}
Alors $T_{\mathbf{R}}$ opère sur $E_{\mathbf{R}}$. Soit $U$ le plus grand
sous-tore compact de $T_{\mathbf{R}}$. Il résulte aisément de l'hyp.\ a)
que
\[
  T_{\mathbf{R}}^{\circ} = U_{\mathbf{R}} \cdot S_{\mathbf{R}} ,
\]
où $S = \mathrm{Im}\,j$, [et $U \subset \mathrm{Ker}\,\varepsilon$].
\struck{Donc} $f$ est de la forme $u\,j(x)$, en d'autres termes il existe
un entier $\nu > 0$ tel que
\[
  f^{\nu} = u \cdot j(x) , \qquad u \in U(\mathbf{R}),\ x \in \mathbf{R}^{*} .
\]
Or $U$ étant compact, les valeurs propres de $u \in U(\mathbf{R})$ sont de
val.\ absolue $1$, donc celles de $f^{\nu}$ sont de val.\ absolue égale à
celles de $j(x)$ ; or les val.\ propres de $j(x)$ \ill{} $x^{i}$ ;
\struck{\ill{} $\varepsilon(u) = 0$} [d'autre part
$\varepsilon | U = \uncertain{0}$], donc
\[
  \varepsilon(f)^{\nu} = \varepsilon(u)\,\varepsilon j(x) = \varepsilon j(x) = x^{\alpha} ,
\]
donc $x^{i/\nu} = \varepsilon(f)^{i/\alpha} = q^{i/\alpha}$, d'où la
conclusion.
\note{« $\varepsilon | U = 0$ » : l'écriture est additive, $\varepsilon$
étant trivial sur $U$}

\struck{N.B. Évidemment \ill{} b) et c) impliquent \ill{} \ill{}
$\mathbf{Q}^{*}$ \ill{}}
\marginal{?}
\note{ce N.B. est biffé d'un long zigzag au crayon ; un « ? » et une
accolade dans la marge gauche}

D'autre part, on voit que a) équivaut au fait que $M$ est, avec les deux
structures supplémentaires provenant de $(j, \varepsilon)$, munissable
d'une structure « positive » des alg.\ motiviques. Ceci, d'autre part, par
un argument \ill{}, implique que
\note{la phrase reste en suspens au bas du feuillet}

\page{93}
$\mathbf{Q}'$ : clôture tot.\ réelle de $\mathbf{Q}$ ….
\note{ligne isolée en tête du feuillet}

\marginal{$\alpha = 2$}
La condition a) équivaut à la suivante, si $q \in \mathbf{Q}^{*}$
\struck{\ill{}} $> 0$ :
\begin{quote}
[Pour tout $\lambda \in M$, $\exists\, r > 0$ tel que]
\struck{Pour tout $\lambda \in M$} $\lambda^{r} \in \mathbf{Q}'(i)^{*}$, et
pour tout $\lambda \in M$, il existe un entier $i$ tel que
$|\lambda| = q^{i/2}$ i.e.\ $\lambda\bar{\lambda} = q^{i}$
\end{quote}
N.B. Le $i$ en question dans \ill{} est déterminé et n'est autre que
$j^{*}(\lambda)$.
\[ \Updownarrow \]
\begin{quote}
Pour tout $\lambda \in M$, \struck{il existe} on a
$|\lambda| = q^{j^{*}(\lambda)/2}$ i.e.\ $\lambda\bar{\lambda} = q^{j^{*}(\lambda)}$
\end{quote}
\note{le « $\alpha = 2$ » de la marge est cerclé ; les deux conditions
sont réunies chacune par une accolade, et la seconde accolade est
seulement ébauchée}

\emph{Cor.} Soit $\lambda \in \overline{\mathbf{Q}}$, tel que
\struck{il existe $q \in \mathbf{Q}^{*+}$, avec}
$\lambda_{\alpha}\bar{\lambda}_{\alpha} = \lambda\bar{\lambda} \in \mathbf{Q}$
($= q$) pour tout conjugué $\lambda_{\alpha}$ de $\lambda$ [i.e.\ $|\lambda_{\alpha}|^{2} = |\lambda|^{2}$]. Alors $\exists\, r > 0$, tel que
$\lambda^{r} \in \mathbf{Q}'(i)$, [dans la condition, on a plongé
\ill{} dans $\overline{\mathbf{Q}} \subset \mathbf{C}$ \ill{}], et \ill{}
plongeant en $\mathbf{Q}(\lambda)$ [relatif à $f$] \ill{} tel que
$\lambda \in M$ (et réciproquement) il existe $M$, $\varepsilon$,
$j$ ….

\marginal{N.B. Ceux des $\lambda$ qui sont réels (donc tot.\ réels !) sont
ceux où $\lambda^{2} = q^{i}$ donc $\lambda = (q^{1/2})^{i}$, les
triviaux}

\marginal{NB. Les conditions n'impliquent pas que
$\lambda \in \mathbf{Q}'(i)^{*}$, comme on voit p.\ ex.\ en prenant
$\lambda = p^{1/4}$ ; Alors $\mathbf{Q}(\lambda)$ n'est ni tot.\ imaginaire ni tot.\ réel ($\mathbf{Q}(p^{1/4})$ est réel), \ill{} tot.\ réel (car il est conjugué à $\mathbf{Q}(\sqrt{-p^{1/2}})$ \ill{}) \ill{}
pas $\subset \mathbf{Q}'(i)$}
\note{les deux N.B. sont dans la marge gauche, en bas du feuillet ; le
second est encadré et découpé par des traits verticaux}

\emph{Proposition.} Soit $\lambda \in \overline{\mathbf{Q}}^{*}$, et
$q \in \mathbf{Q}^{*+}$. Conditions équivalentes :
\begin{enumerate}
  \item[(i)] \struck{$\exists$ entier $i$ tel que pour} [Pour] tout
  conjugué $\lambda_{\alpha}$ de $\lambda$, on ait
  $|\lambda_{\alpha}| = q^{i/2}$ [$\exists$ entier $i$ tel que],
  i.e.\ $\lambda_{\alpha}\bar{\lambda}_{\alpha} = q^{i}$
  \item[(ii)] $\lambda \in \mathbf{Q}'(i)$, il existe un entier $i$ tel
  que $|\lambda| = q^{i/2}$ i.e.\ $\lambda\bar{\lambda} = q^{i}$
  \item[(ii bis)] $\exists$ entier $n > 0$ tel que
  $\lambda^{n} \in \mathbf{Q}'(i)$, et \struck{$\exists$} [entier $i$ tel
  que] $\lambda^{n}\bar{\lambda}^{n} = q^{in}$ ; [$|\lambda| = q^{i/2}$,
  i.e.\ $\lambda\bar{\lambda} = q^{i}$]
  \item[(iii)] $\exists$ un gp de t.m.\ $G$ sur $\mathbf{Q}$, et un
  élément $f$ de $G(\mathbf{Q})$,
  $\mathbf{G}_m \xrightarrow{\;j\;} G \xrightarrow{\;\varepsilon\;} \mathbf{G}_m$,
  avec $\varepsilon j(\lambda) = \lambda^{2}$, \struck{tels que \ill{}}
  [et $\mathrm{Ker}\,\varepsilon$ compact], $\varepsilon(f) = q$, (tels que
  $\lambda$ soit une valeur propre de $f$ [sur $\overline{\mathbf{Q}}$]
  pour un $G$-module de dim.\ finie \ill{}, i.e.\ $\lambda \in$ \ill{} :
  l'image de l'homom.
  $M = D(G_{\overline{\mathbf{Q}}}) \to \overline{\mathbf{Q}}^{*}$
  défini par $f$.
  \item[(iv)] Il existe une algèbre de Riemann commutative
  $A$ [sur $\mathbf{Q}$, et $i \in \mathbf{N}$, et $f \in A$, tels que
  $ff' = q^{i}$] et un homomorphisme $A \to \overline{\mathbf{Q}}$,
  appliquant $f$ dans $\lambda$.
\end{enumerate}
\note{la proposition est remaniée dans l'interligne : le « (ii) »
primitif est surchargé en « (ii bis) », et la condition sur $|\lambda|$
est reprise à droite. L'ajout de (iv) entre crochets est écrit dans la
marge droite, au-dessus de la ligne, et entouré}

\page{95}
\emph{Prop.} Soit $K$ un corps de nbs fini sur $\mathbf{Q}$, $L$ une
extension quadratique de $K$, $\mathfrak{p}$ un idéal premier de $K$,
\struck{$\mathfrak{p}, \mathfrak{p}'$} qui se décompose dans $L$ dans deux
idéaux premiers distincts $\mathfrak{p}$, $\mathfrak{p}'$. Alors il existe
un $x \in L^{*}$ \struck{dont le diviseur soit} \struck{tel que} dont le
diviseur soit $h(\mathfrak{p} - \mathfrak{p}')$ [où $h \in \mathbf{N}^{*}$],
et tel que
\[
  x\,\sigma x = 1 ,
\]
où $\sigma$ est l'automorphisme non trivial de $L/K$.
\marginal{Plus précisément, on trouve que les diviseurs des éléments $x$
de $L$ tels que $N_{L/K}\,x = 1$, sont les diviseurs de $L$ de norme
\uncertain{$1$} !}
\note{note marginale écrite en oblique dans la marge gauche}

En effet, \struck{\ill{}} [soit $h$ l'ordre \ill{} dans le groupe des]
classes de div.\ de $L$, il existe un $y$ de diviseur
$h(\mathfrak{p} - \mathfrak{p}')$. Alors $y\,\sigma y$ est un élément de
$K$ dont le diviseur est évidemment nul, donc c'est une unité $u$.
Posons $z = y^{2}$, alors
\[
  z\,\bar{z} = (y \cdot \sigma y)^{2} = u^{2} ,
\]
[$\mathrm{div}\,x = \mathrm{div}\,z = 2\,\mathrm{div}\,y$] d'où si
$x = z/u = y^{2}/Ny$, on aura
$\mathrm{div}\,x = 2h(\mathfrak{p} - \mathfrak{p}')$, et
\[
  x\,\sigma x = (z\bar{z})/u^{2} = 1 .
\]
\note{« $= 2h(\mathfrak{p} - \mathfrak{p}')$ » est écrit sous
« $x = z/u$ » ; il contredit le « $h(\mathfrak{p} - \mathfrak{p}')$ » de
l'énoncé par un facteur $2$, ce qu'il laisse tel quel}

N.B. Supposons que $K$ est tot.\ réel [et $L$ \struck{\ill{}} ext.\ quadratique \struck{galoisienne de $\mathbf{Q}$}, ($K$, $L$ galoisiens sur
$\mathbf{Q}$)], \struck{et $K$} tot.\ imaginaire. Soit $p$ un nb premier,
et supposons que \ill{} idéal de $K$ \ill{} de $p$, \ill{} premier
\ill{} de $L$ sur $K$. Soit $\lambda$ un \emph{entier} de $L$ tel que
$\lambda\bar{\lambda} = p^{i}$. Alors
$\lambda^{2}\bar{\lambda}^{2} = p^{2i}$ donc si
$\mu = \lambda^{2}/p^{i} = \lambda^{2}/N_{L/K}\lambda$, on aura
\[
  \mu\bar{\mu} = 1 ,
\]
\struck{\ill{}} ce qui implique que $\mu$ est une unité en $\mathfrak{p}$,
donc ($K$ étant galoisien) en \ill{} (l'hyp.\ \ill{} idéal premier \ill{}
$\mathfrak{p}'$), donc $\mu$ est une unité (NB. $\lambda$ \ill{} entier
\ill{}), une racine de l'unité. Donc $\lambda = p^{i/2}u$, $u$ racine de
l'unité !
\note{l'hypothèse sur la décomposition de $p$ dans $L$ n'est pas lue ; la
conclusion demande que les idéaux premiers de $K$ au-dessus de $p$ ne se
décomposent pas dans $L$}

\page{97}
Soit $K$ une extension \struck{\ill{}} tot.\ réelle finie de $\mathbf{Q}$,
$L$ ext.\ quadratique tot.\ imaginaire de $K$, $q = p^{n}$ une puissance
d'un nb premier. On cherche les $\lambda \in K$\note{sic, pour
$\lambda \in L$} entiers tels que $\lambda\bar{\lambda} = q^{i}$, i.e.\ $|\lambda| = q^{i/2}$, [pour un $i$ convenable $\in \mathbf{N}$] (ce sont
des unités exc.\ $p$), et ils engendrent \struck{le groupe des} un
sous-groupe du groupe \struck{\ill{}} des unités exc.\ $p$,
\struck{\ill{}}\struck{\ill{} des $\lambda$ \ill{}} formé des $\lambda$
unités exc.\ $p$ tels que
\[
  \lambda\bar{\lambda} = q^{i} \qquad (i \in \mathbf{Z}) .
\]
Notons que \struck{\ill{}} \struck{\ill{} unités absolues exc.\ $p$. Donc
\ill{} il \ill{} si $q$ \ill{}} \struck{\ill{} $\lambda$ \ill{}
$q^{i/2}u$, avec \ill{}} on peut écrire un $\lambda$ \struck{$K$}
\ill{} comme
\[
  \lambda = q^{i/2}u ,
\]
où $u$ est une unité exc.\ $p$ telle que $u\bar{u} = 1$, ce qui implique
que $u$ est une \emph{unité absolue} exc.\ $p$. Ces $u$ forment un groupe
$U$, et l'homom.
\[
  U \longrightarrow \mathbf{Z}^{P} , \qquad
  u \mapsto (v_{\mathfrak{p}}(u))_{\mathfrak{p} \in P}
\]
où $P$ est l'ensemble des $\mathfrak{p}$ de $L$ au-dessus de $p$
\ill{}, a un noyau fini. D'après la prop.\ plus haut, [son conoyau est
d'indice fini] \ill{} dans $(\mathbf{Z}^{(P)})'$, \ill{} image est
\ill{} formé des diviseurs de la forme
\struck{$\mathfrak{p}' - $ \ill{}} \ill{} $\mathfrak{p}$ \ill{} idéal
premier de \ill{}.
\marginal{(rel.\ à $K$) nulle}
Donc si pour tout $\mathfrak{p}$ \ill{} de $K$ sur $p$ qui
\emph{splitte en deux} dans $L$, on choisit un \ill{} au-dessus, on
trouve un isom.\ \ill{} groupe fini
\[
  U \simeq \mathbf{Z}^{Q} ,
\]
$Q$ étant l'ensemble des \ill{}
\note{le feuillet s'arrête sur ces mots ; la suite du raisonnement n'est
pas au recto suivant, qui reprend les conditions sur les valeurs propres
de Frobenius}

\page{99}
Conditions nécessaires pour qu'un $\lambda \in \overline{\mathbf{Q}}$ soit
valeur propre \struck{\ill{}} de frobenius, pour une variété
\uncertain{projective} \ill{} définie sur le corps $\mathbf{F}_q$ ;
$q = p^{n}$ :

\struck{1) $\lambda \in \mathbf{Q}'(i)$ ; 2) $\lambda^{2}$ est de la forme
$q^{i}u$, où $u \in \mathbf{Q}'(i)$ est une unité exc.\ $p$ ; 3)
$\mathfrak{p}$, $\mathfrak{p}'$ sont deux idéaux premiers de \ill{}}
\note{premier état de la liste, barré}

\begin{enumerate}
  \item[1)] $\lambda$ entier algébrique
  \item[2)] $\lambda \in \mathbf{Q}'(\sqrt{-1})$
  \item[3)] $\lambda^{2} = q^{i}u$, [$i \in \mathbf{Z}$] avec $u$ une unité
  [absolue] exc.\ $p$ (cette dernière condition équivalente à
  [$\exists\, i$ tel que $|\lambda_{\alpha}| = q^{i/2}$ pour tt conjugué
  $\lambda_{\alpha}$ de $\lambda$] \ill{} déterminés \ill{} uniques $i$ et
  $u$)
  \item[4)] \struck{Si $\mathfrak{p}$, $\mathfrak{p}'$ sont deux idéaux
  premiers de $\mathbf{Q}'(\sqrt{-1})$ \ill{} au-dessus d'un idéal premier
  de $\mathbf{Q}'$, on a
  $v_{\mathfrak{p}}(u) + v_{\mathfrak{p}'}(u) = 0$ i.e.\ $v_{\mathfrak{p}}(\lambda) + v_{\mathfrak{p}'}(\lambda) = \frac{in}{2}$,
  $v_{\mathfrak{p}}$ normalisé tel que $v_{\mathfrak{p}}(p) = 1$.}
\end{enumerate}
\marginal{conséquence de 2 et 3 !!!}
\note{les conditions 1) et 2) et la forme équivalente de 3) sont
encadrées, ainsi que la formule
$v_{\mathfrak{p}}(\lambda) + v_{\mathfrak{p}'}(\lambda) = in/2$ ; la
condition 4) est barrée de quatre traits obliques, et une flèche la
relie à la mention cerclée « conséquence de 2 et 3 !!! »}

\marginal{N.B. $q^{1/2} \in \mathbf{Q}'$ !}
\marginal{N.B. Les conditions 2), 3), 4) sont « diophantiennes », la
condition 1) seulement \ill{} On désigne par $M(q^{n})$ le groupe de
$\mathbf{Q}'(i)^{*}$ défini par \struck{3)} 3), 4).}

L'entier $i$ qui intervient ici \ill{} appelé le poids de $\lambda$
(relativement à $q = p^{n}$). \struck{Le poids \ill{}} Si les conditions
\ill{} sont satisfaites, \struck{\ill{}} \ill{} relativement à
$q' = q^{m}$, elles le sont pour $q$, et le poids $i$ rel.\ à $q$ est égal
à $mi'$, où $i'$ est le poids par rapport à $q'$.

\marginal{Th.\ d'existence de motifs sur $\mathbf{F}_q$}
\emph{Conjecture} \struck{\ill{}} \uncertain{fondamentale} : les
\emph{conditions nécessaires précédentes sont suffisantes}.
\note{la conjecture est isolée par un trait vertical ; la note marginale
est marquée d'un double trait vertical}

Cela détermine \struck{\ill{}} donc explicitement, en termes
arithmétiques simples, le groupe de Galois motivique des $\mathbf{F}_q$.
Vérifions que cela donne bien
\note{la phrase se poursuit au feuillet 101 (lot 6)}

\end{document}
